\subsection{Resource Estimation}

The Resource Estimation stage characterizes the resources required by the
execution units identified during the Granularity Analysis stage. It answers a
single question: \emph{how many resources does one instance of each execution
unit need?} Placement and offloadability are deliberately left to the
subsequent stages.

For each execution unit $u_j\in\mathcal{U}$, the LLM analyzes the functions of
its constituent components, the characteristics produced by the Granularity
Analysis (e.g., machine-learning inference, high-frequency input, statefulness,
persistent storage), its unit type (runtime or initialization), and the data
exchanged along its incoming and outgoing edges in the unit graph. When an
estimate depends on information that the functional graph does not provide
(e.g., input resolution or sampling rate), the LLM states the assumption
explicitly, so that every estimate remains traceable.

The estimates are heuristic rather than measured. To keep them traceable, each
resource is derived through an explicit evidence hierarchy: (i) observed facts
stated in the input, (ii) characteristics inferred from these facts, (iii) a
qualitative demand level, (iv) a coarse numerical estimate, and (v) a
confidence level. The qualitative level is the primary result; the numerical
value only refines it and is reported with limited precision, or omitted when
it is not meaningful (e.g., logic executed on a microcontroller). Furthermore,
the demand of the unit's own logic is separated from the baseline overhead of
the runtime platform on which it currently executes (e.g., a mobile
application process); platform baselines are reported once per platform and
are not attributed to individual units.

Computational, memory and storage values are restricted to a fixed estimation
grid (e.g., 0.01, 0.05, 0.1, 0.25, 0.5, 1, 2, 4 cores), which reflects their
order-of-magnitude nature. For platforms whose resources are not comparable
with general-purpose servers, such as microcontrollers, no value in cores or GB
is produced; only the qualitative level is kept, unless the input states a
native figure.
Each profile also carries, unchanged, the constraints identified during the
Granularity Analysis (execution constraint, deployment role, hardware or
platform bindings), so that a low resource demand is never interpreted
independently of the binding that restricts where the unit may execute.
Aggregated demands are computed per platform rather than globally, since the
resources of heterogeneous platforms (microcontrollers, mobile devices,
servers) are not additive. These aggregates describe the load of the
application as it currently executes (application demand plus platform
baseline) and are not capacity requirements for future nodes, since the
Deployment Planner may separate co-located units. Initialization units (e.g.,
model provisioning) are characterized only during their initialization phase:
resources that remain in use afterwards, such as loaded models, are attributed
to the consuming units, which prevents double counting, and their transfers
are described as one-time volumes rather than sustained bandwidth. Finally, every estimate carries a provenance label
\begin{equation}
\pi \in \{\mathrm{observed},\ \mathrm{inferred},\ \mathrm{assumed},\ \mathrm{unknown}\},
\end{equation}
indicating whether the value is stated in the input, derived from stated facts,
dependent on an explicitly listed assumption, or unsupported. When several
facts contribute, the weakest provenance is retained, and an unsupported value
is reported as unknown rather than filled in. Provenance lets the Deployment
Planner distinguish firm constraints from heuristic guesses, e.g., by applying
safety margins to assumed values, and it explains why the precision of the
profile varies between applications with more or less detailed descriptions.
The resulting profile is uncalibrated by construction; its accuracy
is assessed experimentally by comparing the estimates with measurements of the
running application (Section~\ref{sec:evaluation}).

The resource profile of an execution unit is represented as
\begin{equation}
r_j = \left( r_j^{cpu},\; r_j^{mem},\; r_j^{gpu},\; r_j^{sto},\; r_j^{net} \right),
\end{equation}
where $r_j^{cpu}$ is expressed in vCPU cores, $r_j^{mem}$ in GB of RAM,
$r_j^{gpu}$ gives the GPU requirement (none, optional or required) and the
associated GPU memory, $r_j^{sto}$ is the storage in GB, and
$r_j^{net}=(r_j^{in},r_j^{out})$ is the ingress and egress traffic in Mbps.
CPU and memory are estimated both for typical operation and at peak, so that
the Deployment Planner can reason about sustained load and headroom. Each
estimate is also mapped to a qualitative level
(very low, low, medium, high, very high) using fixed thresholds, and is
accompanied by a confidence level.

In addition to per-unit demands, the stage estimates the communication
requirement of every edge $(u_j,u_k)$ of the unit graph as
\begin{equation}
b_{jk} = \frac{s_{jk}\cdot f_{jk}\cdot 8}{1000},
\end{equation}
where $s_{jk}$ is the message size in kB, $f_{jk}$ the message rate per
second, and $b_{jk}$ the required bandwidth in Mbps. Unlike the per-unit
estimates, $b_{jk}$ is neither rounded nor bounded below, because the
Deployment Planner aggregates link loads and rounding small flows would
distort these sums. Each edge also records its end-to-end flow path when
a payload is relayed across several units, and a latency sensitivity defined
on the flow (high: per-input timeliness; medium: seconds-scale; low: eventual
delivery). Each edge is further annotated with its current
communication mechanism and with its boundary type $\beta_{jk}$:
\begin{equation}
\beta_{jk}\in\{\mathrm{logical},\ \mathrm{physical},\ \mathrm{hardware}\}.
\end{equation}
A \emph{logical} edge connects units that currently execute in the same
program or machine (e.g., threads exchanging frames through a queue); its
bandwidth $b_{jk}$ is the network load that would appear only if the planner
separated the two units. A \emph{physical} edge already crosses a link between
devices (e.g., a Bluetooth serial link), and a \emph{hardware} edge denotes a
local hardware interface that cannot be turned into network traffic.
Accordingly, each unit reports both its current physical traffic and its
potential traffic if all its logical edges became network links.

The complete resource profile of the application is
\begin{equation}
R = \left( \{r_1,\ldots,r_m\},\; \{b_{jk}\},\; \Pi \right),
\end{equation}
where $\Pi$ denotes the set of platform baselines. All values refer to a single instance of an execution unit and use the same
units as the capacity model of the infrastructure, which allows the
Offloadability Analysis and the Deployment Planner to compare demands with
node capacities directly.

The procedure is summarized in Algorithm~\ref{alg:resource_estimation}.

\begin{algorithm}[ht]
\caption{LLM-based Resource Estimation}
\label{alg:resource_estimation}
\begin{algorithmic}[1]
\Require Functional application graph $G=(V,E)$
\Require Execution units $\mathcal{U}$ and unit graph $G_U=(\mathcal{U},E_U)$
\Ensure Resource profile $R$
\State $R_U \leftarrow \emptyset$, $B \leftarrow \emptyset$
\For{each execution unit $u_j \in \mathcal{U}$}
  \State $I_j \leftarrow \mathrm{GetUnitInformation}(u_j, G, G_U)$
  \State $r_j \leftarrow \mathrm{LLMResourceEstimation}(u_j, I_j)$
  \State $R_U \leftarrow R_U \cup \{r_j\}$
\EndFor
\For{each edge $(u_j,u_k) \in E_U$}
  \State $b_{jk} \leftarrow \mathrm{LLMCommunicationEstimation}(u_j,u_k,G_U)$
  \State $B \leftarrow B \cup \{b_{jk}\}$
\EndFor
\State $R \leftarrow (R_U, B)$
\State \Return $R$
\end{algorithmic}
\end{algorithm}
